\section{Математическое описание}

\subsection{Генерация графа и основные понятия}

Графом $G(V, E)$ называется совокупность двух множеств:
\begin{itemize}
    \item непустое множество $V$~--- множество вершин;
    \item $E$, множество двуэлементных подмножеств множества $V$,~--- множество рёбер.
\end{itemize}

Ребро можно трактовать не только как множество ${\{v_1, v_2\}}$, но и как пару вершин $(v_1, v_2)$.

\[
p \stackrel{\text{Def}}{=} p(G) \stackrel{\text{Def}}{=} |V|, \hspace{2em}
q \stackrel{\text{Def}}{=} q(G) \stackrel{\text{Def}}{=} |E|.
\]

Если рёбра задаются упорядоченными парами, то граф называется ориентированным. В противном случае --- неориентированным. Степенью вершины $d(v)$ называется количество рёбер, инцидентных этой вершине.

Для неориентированного графа справедлива лемма о рукопожатиях. Сумма степеней вершин графа равна удвоенному количеству рёбер.

\[
\sum_{v \in V} d(v) = 2q.
\]

Маршрутом в графе называется чередующаяся последовательность вершин и рёбер, начинающаяся и заканчивающаяся вершиной, в которой любые два соседних элемента инцидентны.

\begin{itemize}
    \item Если все рёбра различны, то маршрут называется цепью.
    \item Если все вершины различны, то маршрут называется простой цепью.
\end{itemize}

Если $v_0 = V_k$, то маршрут замкнут, иначе~--- открыт.

\vspace{1em}

Цепь, соединяющая вершины $u$ и $v$, обозначается $\langle u, v\rangle$. Говорят, что две вершины в графе связны, если существует соединяющая их цепь. \textit{Граф, в котором все вершины связны, называется связным}.

Замкнутая цепь называется циклом, замкнутая простая цепь называется простым циклом. Число циклов в графе обозначается $z(G)$. \textit{Граф без циклов называется ациклическим}.


\vspace{1em}

В первой лабораторной работе связный ациклический граф генерируется случайным образом на основе смещённого распределения Вейбулла. Распределение задаётся функцией, приведённой в справочнике Р.~Н.~Вадзинского:

\[
F(y) = 1 - \exp\left(-\left(\frac{y-y_0}{a}\right)^c\right),
\]

где $y$~--- случайная величина, $a$~--- параметр масштаба, $c$~--- параметр формы $(c > 0)$, $y_0$~--- параметр смещения.

Для генерации случайных чисел по смещённому распределению Вейбулла используется формула
\[
y = y_0 + a(-\ln r)^{1/c}.
\]
Здесь $r$~--- равномерно распределённая на интервале $(0, 1)$ случайная величина.

На рис. \ref{fig:weibull-shape-comparison} приведены гистограммы 1000 сгенерированных значений при $a=1$, $c=0{,}5; 1{,}5; 2{,}5$ и $y_0=0$. При $c<1$ большая часть значений сосредоточена возле нуля. При увеличении $c$ значения становятся более сконцентрированными, а максимум гистограммы смещается вправо.

\img[1]{tex/graphics/python/weibull-shape-comparison.png}{Гистограммы смещённого распределения Вейбулла при различных значениях $c$}{\label{fig:weibull-shape-comparison}}

Для генерации матрицы смежности использовались параметры $a=0.17$, $c=2.0$, $y_0=0.05$. При формировании графа значение степени $i$-й вершины вычисляется по формуле

\[
d_i = \operatorname{round}\left((p-1)y_i\right),
\qquad i=1,\ldots,p,
\]

где $p$~--- количество вершин графа, $y_i$~--- значение случайной величины, полученное по смещённому распределению Вейбулла. Для графа из 10 вершин формула принимает вид

\[
d_i = \operatorname{round}(9y_i).
\]

Граф строится по полученной последовательности степеней вершин при выполнении условия

\[
1 \leq d_i \leq p-1
\qquad \text{для всех } i=1,\ldots,p.
\]

Если хотя бы одна степень не удовлетворяет этому условию, последовательность генерируется заново. Такой набор параметров подобран экспериментально. Он обеспечивает умеренные степени вершин и не приводит к формированию почти полного графа.

Гистограмма 1000 значений распределения с выбранными параметрами приведена на рис. \ref{fig:weibull-graph-generation}. На том же рисунке показаны рассчитанные значения степеней для графа из 10 вершин.

\img[1]{tex/graphics/python/weibull-graph-generation.png}{Распределение Вейбулла и значения степеней при $p=10$}{\label{fig:weibull-graph-generation}}

В рамках работы смещённое распределение Вейбулла используется также для генерации весов рёбер и сети потоков. Для весовой матрицы пользователь может выбрать знак генерируемых весов: положительные, отрицательные или смешанные. При генерации со смешанными весами знак выбирается случайно.

Расстоянием между вершинами $u$ и $v$ (обозначается $d(u, v)$) называется длина кратчайшей цепи $\langle u, v\rangle$, а сама кратчайшая цепь называется геодезической.

\[
d(u, v) \stackrel{\text{Def}}{=} \min_{\langle u, v\rangle} \langle u, v\rangle.
\]

\begin{center}
Если $\nexists (\langle u, v\rangle)$, то $d(u, v) \stackrel{\text{Def}}{=} +\infty$.
\end{center}

Диаметром графа $G$ называется длиннейшая геодезическая. Длина диаметра обозначается $D(G)$:

\[
D(G) \stackrel{\text{Def}}{=} \max_{u, v \in V} d(u, v).
\]

Эксцентриситетом $e(v)$ вершины $v$ в связном графе $G(V, E)$ называется максимальное расстояние от $v$ до любой другой вершины графа $G$:

\[
e(v) \stackrel{\text{Def}}{=} \max_{u \in V} d(u, v).
\]

Радиусом $R(G)$ графа $G$ называется наименьший из эксцентриситетов вершин:

\[
R(G) \stackrel{\text{Def}}{=} \min_{v \in V} e(v).
\]

Вершина $v$ называется центральной, если её эксцентриситет совпадает с радиусом графа, т.е. $e(v) = R(G)$.

Множество центральных вершин называется центром графа и обозначается $C(G)$:

\[
C(G) \stackrel{\text{Def}}{=} \{v \in V \mid e(v) = R(G)\}.
\]

Метод Шимбелла использует степени весовой матрицы для нахождения весов маршрутов с фиксированным количеством рёбер. В данной работе для каждой пары вершин определяется минимальный или максимальный суммарный вес маршрута, содержащего ровно $k$ рёбер.

Пусть $W=(w_{ij})$~--- весовая матрица графа, а $\omega_{ij}^{(k)}$~--- минимальный или максимальный вес маршрута из вершины $v_i$ в вершину $v_j$, содержащего ровно $k$ рёбер. При Шимбелловом умножении произведение элементов заменяется сложением весов, а сложение полученных значений~--- выбором минимума или максимума. Правила выполнения операций приведены на рис.~\ref{fig:shimbell}.

\img[1][trim=0 75 0 0,clip]{img/shimbell.png}{Матричные операции метода Шимбелла}{\label{fig:shimbell}}

При возведении весовой матрицы во вторую степень элемент результирующей матрицы вычисляется по формуле

\[
\omega_{ij}^{(2)} = \min(\max)\left\{
    \left(\omega_{i1}^{(1)} + \omega_{1j}^{(1)}\right),
    \left(\omega_{i2}^{(1)} + \omega_{2j}^{(1)}\right),
    \ldots,
    \left(\omega_{ip}^{(1)} + \omega_{pj}^{(1)}\right)
\right\}.
\]

Для произвольного количества рёбер вычисления продолжаются рекуррентно. Произведение матриц в следующей формуле является Шимбелловым умножением и выполняется по правилам, приведённым на рисунке \ref{fig:shimbell}:

\[
W^{(1)}=W,
\qquad
W^{(k)}=W^{(k-1)} \cdot W.
\]

При $k=0$ диагональные элементы матрицы равны нулю, а остальные элементы соответствуют отсутствующим маршрутам.

\subsection{Связность и точки сочленения}

Отношение связности вершин называется эквивалентностью. \textit{Классы эквивалентности по отношению связности называются компонентами связности графа}.

Вершина графа называется точкой сочленения, если её удаление увеличивает число компонент связности. Мостом называется ребро, удаление которого увеличивает число компонент связности.

Для поиска точек сочленения граф обходится в глубину. Для каждой вершины определяется, существует ли из её поддерева путь к ранее посещённым вершинам в обход неё. Если такого пути нет, удаление этой вершины разделит граф, поэтому она является точкой сочленения. Для корневой вершины отдельно проверяется количество независимых ветвей обхода. В реализации ориентированный граф рассматривается как неориентированный.

Псевдокод обхода в глубину с вычислением нижних меток приведён на рис.~\ref{fig:articulation-points}.

\img[0.75]{img/atricupation-points.png}{Псевдокод поиска точек сочленения}{\label{fig:articulation-points}}

Стандартный алгоритм Дейкстры находит кратчайшие пути от начальной вершины $s$ до остальных вершин графа с неотрицательными весами. На каждой итерации выбирается вершина с минимальным текущим весом пути. Затем веса путей до соседних вершин пересчитываются в соответствии с неравенством треугольника. После этого выбранная вершина считается обработанной окончательно. Псевдокод стандартного алгоритма приведён на рис.~\ref{fig:dijkstra-positive}.

\img[0.7]{img/dijkstra-positive.png}{Псевдокод стандартного алгоритма Дейкстры}{\label{fig:dijkstra-positive}}

Для работы с отрицательными весами алгоритм модифицирован: вершина не фиксируется окончательно и при нахождении пути меньшего веса повторно помещается в очередь с приоритетами. Поэтому её вес может уменьшаться несколько раз. Псевдокод модифицированного алгоритма приведён на рис.~\ref{fig:dijkstra-negative}.

\img[0.5][trim=324 10 30 259,clip]{img/dijkstra-negative.png}{Псевдокод алгоритма Дейкстры для графов с отрицательными весами}{\label{fig:dijkstra-negative}}

Для защиты от зацикливания алгоритм подсчитывает количество уменьшений веса каждой вершины. Если для некоторой вершины оно достигает количества вершин графа $p$, выполнение прекращается и фиксируется наличие отрицательного цикла. При его отсутствии алгоритм завершается после опустошения очереди.

Результаты сравнения алгоритмов по количеству итераций и соответствующие графики приведены в \hyperref[sec:work-results]{разделе «Результаты работы»}. По результатам запусков поиск точек сочленения оказался эффективнее модифицированного алгоритма Дейкстры.

\subsection{Максимальный поток и поток минимальной стоимости}

Пусть $G = (V, E)$~--- орграф с одним истоком $s \in V$ и одним стоком $t \in V$, где $s \neq t$.

Сетью $S = (G; c)$ называется орграф $G$ с заданной функцией $c : E \rightarrow \mathbb{R}_{+}$, сопоставляющей каждой дуге $e \in E$ неотрицательное число $c(e)$, называющееся пропускной способностью этой дуги $e$.

Потоком в сети $S = (G; c)$, где $G = (V, E)$, называется функция $f: E \rightarrow \mathbb{R}_{+}$, удовлетворяющая условиям ограничения пропускной способности

\[
0 \leq f(u,v) \leq c(u,v), \qquad (u,v)\in E,
\]

и сохранения потока в каждой вершине, кроме истока и стока:

\[
\sum_{(u,v)\in E} f(u,v)
=
\sum_{(v,w)\in E} f(v,w),
\qquad v\in V\setminus\{s,t\}.
\]

Величиной потока называется суммарный поток, выходящий из истока, за вычетом потока, входящего в него:

\[
F(f)=\sum_{(s,v)\in E}f(s,v)-\sum_{(v,s)\in E}f(v,s).
\]

Разрезом сети относительно истока $s$ и стока $t$, или $(s,t)$-разрезом, называется разбиение множества вершин $V$ на два непересекающихся подмножества $V_s$ и $V_t$ такие, что $V_s\cup V_t=V$, $s\in V_s$ и $t\in V_t$. Обозначим этот разрез через $P=(V_s,V_t)$. Пропускной способностью разреза $C(P)$ называется суммарная пропускная способность дуг, идущих из $V_s$ в $V_t$:

\[
C(P)=\sum_{\substack{(u,v)\in E\\u\in V_s,\,v\in V_t}}c(u,v).
\]

\begin{samepage}
Теорема Форда~--- Фалкерсона утверждает, что максимальная величина потока в сети равна минимальной пропускной способности среди всех $(s,t)$-разрезов этой сети:

\[
\max_f F(f)=\min_P C(P).
\]
\end{samepage}

Для поиска максимального потока применяется алгоритм Форда~--- Фалкерсона. В данной работе он реализован с помощью расстановки меток. Сначала поток по всем дугам полагается равным нулю, а исток получает начальную метку. Далее метки распространяются по дугам, по которым поток можно увеличить, и по обратным дугам, по которым ранее проведённый поток можно уменьшить. Каждая метка хранит предыдущую вершину, направление изменения потока и наименьшую пропускную способность найденной цепи.

Если метка достигает стока, по сохранённым вершинам восстанавливается увеличивающая цепь и поток вдоль неё изменяется на найденную величину. После этого метки расставляются заново. Когда сток больше нельзя пометить, увеличивающих цепей не остаётся и полученный поток является максимальным. Псевдокод алгоритма приведён на рис.~\ref{fig:max-flow-algorithm}.

\imgGrid[0pt](X[c,m]){
    \gridImage[width=0.6\linewidth,trim=0 41 40 0,clip]{img/max-flow-1.png} \\
    \gridImage[width=0.52\linewidth]{img/max-flow-2.png}
}{Псевдокод алгоритма Форда~--- Фалкерсона}{\label{fig:max-flow-algorithm}}

Рассмотрим задачу определения потока заданной величины $\theta$ от истока $s$ к стоку $t$ в сети

\[
G=(S,U,\Omega,D),
\]

где $S$~--- множество вершин, $U$~--- множество дуг, $\Omega$~--- множество их пропускных способностей, $D$~--- множество стоимостей прохождения единицы потока. Каждой дуге $(x_i,x_j)\in U$ соответствуют пропускная способность $c(x_i,x_j)\in\Omega$ и неотрицательная стоимость $d(x_i,x_j)\in D$. Если $\theta>\varphi_{\max}$, где $\varphi_{\max}$~--- величина максимального потока из $s$ в $t$, решения не существует. При $\theta\leq\varphi_{\max}$ допустимых потоков может быть несколько; среди них требуется найти поток минимальной стоимости.

Математическая модель задачи имеет вид

\[
\min_{\varphi}\sum_{(x_i,x_j)\in U}
d(x_i,x_j)\varphi(x_i,x_j)
\]

при следующих ограничениях:

\begin{enumerate}
    \item для каждой дуги $(x_i,x_j)\in U$
    \[
    0\leq\varphi(x_i,x_j)\leq c(x_i,x_j);
    \]

    \item для истока $s$
    \[
    \sum_{x_j\in S}\varphi(s,x_j)-
    \sum_{x_j\in S}\varphi(x_j,s)=\theta;
    \]

    \item для стока $t$
    \[
    \sum_{x_j\in S}\varphi(t,x_j)-
    \sum_{x_j\in S}\varphi(x_j,t)=-\theta;
    \]

    \item для каждой вершины $x_i\in S\setminus\{s,t\}$
    \[
    \sum_{x_j\in S}\varphi(x_i,x_j)-
    \sum_{x_j\in S}\varphi(x_j,x_i)=0.
    \]
\end{enumerate}

Для потока $\varphi$ исходная сеть преобразуется в модифицированную сеть

\[
\widetilde G=(\widetilde S,\widetilde U,\widetilde\Omega,\widetilde D),
\]

которая строится следующим образом:

\begin{enumerate}
    \item $\widetilde S=S$, то есть множество вершин не изменяется;

    \item $\widetilde U=U\cup V$, где $V$~--- множество фиктивных дуг;

    \item если $(x_j,x_i)\in U$ и $\varphi(x_j,x_i)>0$, то фиктивная дуга $(x_i,x_j)$ включается в множество $V$. Для неё полагается
    \[
    \widetilde c(x_i,x_j)=\varphi(x_j,x_i),
    \qquad
    \widetilde d(x_i,x_j)=-d(x_j,x_i);
    \]

    \item для ненасыщенной дуги $(x_i,x_j)\in U$, для которой $\varphi(x_i,x_j)<c(x_i,x_j)$ и $\varphi(x_j,x_i)=0$, полагается
    \[
    \widetilde c(x_i,x_j)=c(x_i,x_j)-\varphi(x_i,x_j),
    \qquad
    \widetilde d(x_i,x_j)=d(x_i,x_j);
    \]

    \item для насыщенной дуги $(x_i,x_j)\in U$, для которой $\varphi(x_i,x_j)=c(x_i,x_j)$ и $\varphi(x_j,x_i)=0$, полагается
    \[
    \widetilde c(x_i,x_j)=0,
    \qquad
    \widetilde d(x_i,x_j)=\infty.
    \]
\end{enumerate}

Множества $\widetilde\Omega$ и $\widetilde D$ образуются соответственно из значений $\widetilde c(x_i,x_j)$ и $\widetilde d(x_i,x_j)$ для всех дуг $(x_i,x_j)\in\widetilde U$.

\subsection{Остовы, код Прюфера и независимые множества рёбер}

Граф без циклов называется ациклическим, или лесом. Связный ациклический граф называется свободным деревом.

Пусть $G = (V, E)$~--- граф. Остовным подграфом графа $G$ называется подграф, содержащий все вершины графа $G$. Остовный подграф, являющийся деревом, называется остовом. Поскольку вершины остова совпадают с вершинами исходного графа, остов однозначно задаётся множеством своих рёбер. Несвязный граф не имеет остова, а связный граф может иметь несколько остовов.

Матрицей Кирхгофа графа $G$ называется матрица $B(G)=(\beta_{ij})$, определяемая как разность диагональной матрицы степеней вершин и матрицы смежности:

\[
B(G)=
\begin{pmatrix}
\deg v_1 & 0 & \cdots & 0 \\
0 & \deg v_2 & \cdots & 0 \\
\vdots & \vdots & \ddots & \vdots \\
0 & 0 & \cdots & \deg v_p
\end{pmatrix}
-A(G).
\]

Иными словами,

\[
\beta_{ij}=
\begin{cases}
-1, & i\neq j \text{ и вершины } v_i \text{ и } v_j \text{ смежны},\\
0, & i\neq j \text{ и вершины } v_i \text{ и } v_j \text{ не смежны},\\
\deg v_i, & i=j.
\end{cases}
\]

\textit{Теорема Кирхгофа.} Число остовов связного графа $G$ порядка $p\geq2$ равно алгебраическому дополнению любого элемента матрицы Кирхгофа $B(G)$. Все алгебраические дополнения матрицы $B(G)$ равны между собой.

Минимальным остовом взвешенного графа называется остов с наименьшей суммой весов рёбер. Алгоритм Краскала рассматривает рёбра в порядке возрастания их весов. В начале каждая вершина образует отдельное множество. Очередное ребро добавляется в остов, если его концы принадлежат разным множествам, после чего эти множества объединяются. Такая проверка не допускает образования циклов. После выбора $|V|-1$ ребра получается минимальный остов. Псевдокод алгоритма приведён на рис.~\ref{fig:kruskal}.

\img[0.4][trim=1 1 1 1,clip]{img/kruskal-algorithm.png}{Псевдокод алгоритма Краскала}{\label{fig:kruskal}}

Множество рёбер называется независимым, или паросочетанием, если никакие два его ребра не смежны. Наибольшее число рёбер в независимом множестве рёбер называется рёберным числом независимости и обозначается $\beta_{1}$.

Вершина, не инцидентная ни одному ребру паросочетания, называется свободной. Увеличивающий путь начинается и заканчивается в свободных вершинах, а его рёбра попеременно не принадлежат и принадлежат паросочетанию. Замена входящих в паросочетание рёбер такого пути на не входящие увеличивает число выбранных рёбер на один.

Алгоритм Эдмондса ищет увеличивающие пути обходом в ширину. Если в дереве поиска обнаруживается цикл нечётной длины, он временно сжимается в одну вершину, после чего поиск продолжается. При нахождении увеличивающего пути сжатые циклы раскрываются, а паросочетание увеличивается. Если увеличивающих путей больше нет, полученное паросочетание является максимальным. Псевдокод алгоритма приведён на рис.~\ref{fig:edmonds-blossom}.

Подробное описание и интерактивная визуализация алгоритма приведены в материале Технического университета Мюнхена~\cite{tum-blossom}.

\img[0.55][trim=10 9 10 110,clip]{img/edmonds-blossom-algorithm.png}{Псевдокод алгоритма Эдмондса}{\label{fig:edmonds-blossom}}

\subsection{Эйлеровы циклы и фундаментальная система разрезов}

Если граф имеет цикл (не обязательно простой), содержащий все рёбра графа, то такой цикл называется эйлеровым циклом, а граф~--- эйлеровым графом. \textit{Эйлеровым может быть только связный граф.}

Для приведения графа к эйлерову виду сначала находятся вершины нечётной степени. Из них выбираются две вершины $u$ и $v$. Если они не смежны, ребро $(u,v)$ добавляется. Если ребро существует, оно удаляется только при наличии обходного пути между $u$ и $v$, что сохраняет связность графа. После каждой операции чётность степеней $u$ и $v$ меняется, поэтому число вершин нечётной степени уменьшается на два. Поскольку таких вершин всегда чётное число, процесс завершается за конечное число шагов.

Для построения эйлерова цикла применяется алгоритм Флери. Начальная вершина помещается в стек. Из текущей вершины выбирается соседняя вершина, после чего соединяющее их ребро удаляется, а выбранная вершина добавляется в стек. Переходы продолжаются, пока у текущей вершины не останется соседей. Такая вершина добавляется в результат и удаляется из стека, после чего обход продолжается из предыдущей вершины. После опустошения стека получается эйлеров цикл. Материал по построению эйлеровых циклов приведён на сайте Технического университета Мюнхена~\cite{tum-eulerian-cycle}. Псевдокод алгоритма приведён на рис.~\ref{fig:eulerian-cycle-algorithm}.

\img[0.9][trim=0 53 0 0,clip]{img/fluery-algorithm.png}{Псевдокод алгоритма Флери}{\label{fig:eulerian-cycle-algorithm}}

Пусть $T(V, E_{T})$~--- некоторый остов графа $G(V, E)$. Рассмотрим ребро остова $e$ и определим разрез $S_{e}$ слеующим образом.

Ребро $e$~--- мост в дереве $T$. Следовательно, удаление ребра $e$ разбивает множество вершин $V$ на два непустых подмножества $V_{1}$ и $V_{2}$, так что

\[
V_1 \subset V, V_1 \neq \varnothing, V_2 \subset V, V_2 \neq \varnothing, V_1 \cup V_2 = V, V_1 \cap V_2 = \varnothing.
\]

Включим в разрезе $S_{e}$ ребро $e$ и все хорды остова $T$, соединяющие вершины множества $V_{1}$ с вершинами множества $V_{2}$, тогда $S_{e}$~--- простой разрез:

\[
S_{e} \stackrel{\text{Def}}{=} e + \{(v_1, v_2) \in T \mid v_1 \in V_1~\&~v_2 \in V_2\} = E(V_1, V_2).
\]

Система разрезов называется фундаментальной системой разрезов. Разрезы фундаментальной системы называются фундаментальными, а количесово разрезов в фунламентальной системе называется коциклическим рангом графа $G$ и обозначается $m^{*}(G)$.

Рассмотрим векторное пространство над двоичной арифметикой, натянутое на множество рёбер $E$ графа $G(V, E)$. Элемент этого векторного пространства (линейную комбинацию) можно отождествлять с подмиодеством рёбер: если коэффициент в линейной комбинации равен $1$, то ребро входит в подмножество, если коэффициент равен $0$, то не входит.

Сложению векторов соответствует симметричесая разность множeств рёбер. Множество правильных разрезов замкнуто относительно симметрической разности, поэтому оно является подпространством общего пространства рёбер графа.

Множество разрезов называется независимым, если ни один из разрезов не является линейной комбинацией остальных.

Максимальное независимое множество разрезов является базисом подпространства разрезов. Фундаментальная система разрезов является базисом подпространства разрезов. Разрезы любой фундаментальной системы незаыисимы.

Для построения фундаментальной системы сначала алгоритмом Краскала строится остов $T$. Затем из остова исключается ребро $e$, в результате чего остов разделяется на две части. В разрез $S_e$ включаются исключённое ребро и все хорды, соединяющие вершины разных частей. Эти действия повторяются для каждого ребра остова, то есть $p-1$ раз. Полученные разрезы образуют фундаментальную систему.
